\chapter{Introduction}\label{ch:intro}

\section{What PlanWise does}

PlanWise is a scheduling agent for a make-to-order plant. It takes the orders customers have placed, the
plant's machines, shifts and rules, and the stock that is already on the floor, and it builds a
\term{reasonable} production schedule for the next planning period: one that is feasible on every machine,
delivers as many orders on time as possible, keeps finished goods above their safety level, and does not open
more shifts or changeovers than necessary.

The planner stays in charge. PlanWise proposes several schedules (\term{options}), explains how each was
built and why an order is late, lets the planner edit any of them by hand, and records which one was
committed and why. The committed schedule then becomes printable work orders for the shop floor, and can be
replanned in minutes when a machine breaks down.

The example used throughout this guide is an alloy wheel plant with five stages:

\begin{center}
\begin{tikzpicture}[stage/.style={draw=brand, fill=brandsoft, rounded corners=3pt, minimum width=2.45cm, minimum height=1.1cm, align=center, font=\small\bfseries}, arr/.style={-{Stealth[length=2.4mm]}, thick, brand}]
\node[stage] (c) {Casting\\[-1pt]{\scriptsize\normalfont 2 die-casting cells}};
\node[stage, right=0.55cm of c] (h) {Heat treatment\\[-1pt]{\scriptsize\normalfont 2 T6 furnaces, 24/7}};
\node[stage, right=0.55cm of h] (m) {Machining\\[-1pt]{\scriptsize\normalfont 2--3 CNC lathes}};
\node[stage, right=0.55cm of m] (p) {Painting\\[-1pt]{\scriptsize\normalfont conveyor line}};
\node[stage, right=0.55cm of p] (q) {Test \& packing\\[-1pt]{\scriptsize\normalfont 1 station}};
\draw[arr] (c) -- (h); \draw[arr] (h) -- (m); \draw[arr] (m) -- (p); \draw[arr] (p) -- (q);
\end{tikzpicture}
\end{center}

Four products (17-inch front and rear wheels, each in silver and matte black) pass through all five stages.
Each stage works in fixed-size \term{runs} (for example 16 wheels per casting run, 24 per furnace load). Front
and rear wheels share their semi-finished stock up to machining; the colour is only decided at the paint line.
Heat treatment takes eight hours per load and is the plant's bottleneck.

\section{Who this guide is for}

\begin{itemize}
  \item \textbf{Production planners} prepare each period, generate and compare schedules, edit and commit
        one, and replan after incidents. Chapters~\ref{ch:today}--\ref{ch:variants} are written for them.
  \item \textbf{Plant or industrial engineers} maintain the factory set-up: stages, products, machines, molds,
        shifts and the policies that define a good schedule (Chapter~\ref{ch:factory}).
  \item \textbf{Shift supervisors} mainly use the work orders of the committed schedule
        (Section~\ref{sec:work-orders}).
  \item \textbf{Analysts and researchers} use the analysis views and the Advanced section to compare
        algorithms on the same data (Chapters~\ref{ch:analysis} and~\ref{ch:advanced}).
\end{itemize}

\section{Conventions used in this guide}

\begin{itemize}
  \item Labels exactly as they appear on screen are set in bold sans-serif, for example \ui{Generate options}
        or \ui{Save \& check}.
  \item Keys are shown as \kbd{Ctrl}\,\kbd{K}.
  \item Codes and identifiers are set in monospace, for example \code{F\_SILVER} or \code{CNC\_1}.
  \item Numbered red markers such as \callout{1} \callout{2} \callout{3} on a screenshot are explained in the
        list directly below that figure.
  \item Boxes mark tips \faLightbulb[regular], notes \faInfoCircle{} and cautions \faExclamationTriangle.
\end{itemize}

\section{Key concepts}\label{sec:concepts}

PlanWise separates data by \emph{how often it changes}. Getting this idea right makes every screen easier to
understand.

\begin{center}
\begin{tabularx}{\linewidth}{>{\bfseries}p{3.1cm}Xp{3.2cm}}
\toprule
\rowcolor{white}Layer & What it contains & Changes \\
\midrule
Factory set-up & Stages, run sizes, products and routes, machines with rates and changeovers, molds, shifts, rest days, holidays, weekly maintenance, objective policies. Every save creates an immutable \term{version}. & A few times a year (Factory page) \\
Order book & Every customer order with product, quantity, due date, priority and customer. & Continuously (Order book page) \\
Period inputs & For one planning period: the orders taken from the book, opening stock, mold counters, calendar exceptions, policy and pull-ahead choice. & Every period, usually two weeks (Prepare step) \\
Derived data & Lots, runs per stage, allocation of stock to orders, checkpoints, machine availability. & Never typed in: computed when the period is saved \\
\bottomrule
\end{tabularx}
\end{center}

\subsection*{Glossary of the most important terms}

\begin{description}[style=nextline, leftmargin=1.2em, font=\bfseries\color{brand}]
  \item[Planning period] A window (by default 14 days) that is scheduled as one problem. It takes the factory
        set-up version valid on its start date and the open orders due up to its end.
  \item[Option / schedule] One complete, feasible schedule of a period produced by an algorithm or by a
        planner's edit. A period usually has several options to compare.
  \item[Run (solver run)] One execution of one algorithm on one copy of the data. Each option is the result of
        a run; every run keeps its own data snapshot so it can be reproduced.
  \item[Run (production run)] The unit of work at a stage: a fixed quantity of one item processed on one
        machine, for example ``F\_CAST run 07, 16 units on CAST\_1''.
  \item[Semi-finished code] The identity of goods between stages (for example \code{F\_HEAT}: front wheel,
        heat-treated). Products that share a code share that stock.
  \item[Safety stock] The minimum finished-goods level a product should keep. Falling below it at an
        end-of-day checkpoint counts as \term{safety-stock shortfall}.
  \item[Stock cap] The maximum finished-goods stock the warehouse can hold for a product.
  \item[Policy] The rule that says which schedule is better: a weighted score over objective terms (lateness,
        shortfall, setup time and so on), optionally with ``strictly first'' terms compared before the others.
  \item[Commit] The planner's decision that one option is the plan. It is logged with a mandatory reason.
  \item[Variant] A what-if copy of a period's inputs, used to test a change without touching the main plan.
  \item[Reschedule] A version of a period created after an incident: everything before the \term{cut-off}
        is frozen, the rest is planned again.
  \item[Pull-ahead (build-ahead)] Using spare capacity to make part of the orders due after the period early.
  \item[Look-ahead window] The days after the period whose orders are shown as ``after the period'' and may
        be pulled ahead (28 days by default).
\end{description}

\section{The planning cycle}\label{sec:cycle}

A planning period goes through the same steps every time. The period page shows them as tabs, and the
status badge tells you which step comes next.

\begin{center}
\begin{tikzpicture}[node distance=0.42cm, box/.style={draw=brand, fill=brandsoft, rounded corners=3pt, minimum height=1.05cm, text width=2.35cm, align=center, font=\small}, arr/.style={-{Stealth[length=2.2mm]}, thick, brand}]
\node[box, fill=white] (book) {\textbf{Order book}\\{\scriptsize orders entered once}};
\node[box, right=of book] (create) {\textbf{Create period}\\{\scriptsize dates only}};
\node[box, right=of create] (prep) {\textbf{1 Prepare}\\{\scriptsize orders, stock, calendar}};
\node[box, right=of prep] (gen) {\textbf{2 Options}\\{\scriptsize generate \& rank}};
\node[box, below=0.9cm of gen] (cmp) {\textbf{3 Schedule \& compare}\\{\scriptsize Gantt, why late, edit}};
\node[box, left=of cmp] (commit) {\textbf{4 Commit}\\{\scriptsize with a reason}};
\node[box, left=of commit] (exec) {\textbf{5 Execute}\\{\scriptsize work orders, incidents}};
\node[box, left=of exec, fill=white] (close) {\textbf{Close \& roll over}\\{\scriptsize next period}};
\draw[arr] (book) -- (create); \draw[arr] (create) -- (prep); \draw[arr] (prep) -- (gen); \draw[arr] (gen) -- (cmp);
\draw[arr] (cmp) -- (commit); \draw[arr] (commit) -- (exec); \draw[arr] (exec) -- (close);
\draw[arr, dashed] (close.north) -- (create.south) node[midway, left, font=\scriptsize, text=muted]{closing stock = opening stock};
\draw[arr, dashed] (exec.north) to[bend left=20] node[midway, right, font=\scriptsize, text=muted, align=left]{incident:\\reschedule} (prep.south);
\end{tikzpicture}
\end{center}

\begin{enumerate}
  \item \textbf{Order book.} Sales or the planner enter orders once, by hand or by importing a spreadsheet
        (Chapter~\ref{ch:orders}).
  \item \textbf{Create the period.} Pick a start date. The set-up, shifts, rest days, open orders and the
        projected opening stock come in automatically (Section~\ref{sec:new-period}).
  \item \textbf{Prepare.} Review the orders, correct the opening stock to the physical count, add one-off
        machine stops, pick the policy, and save so the inputs are checked (Chapter~\ref{ch:prepare}).
  \item \textbf{Generate options.} Choose how long to wait; PlanWise runs suitable algorithms in the
        background and recommends the schedules worth considering (Chapters~\ref{ch:generate}
        and~\ref{ch:options}).
  \item \textbf{Schedule \& compare.} Inspect the Gantt chart, compare two options, see why an order is late,
        and edit a schedule by hand if needed (Chapters~\ref{ch:schedule} and~\ref{ch:editing}).
  \item \textbf{Commit.} Choose one option and record why (Section~\ref{sec:commit}).
  \item \textbf{Execute.} Print work orders, replan after incidents, and finally close the period. Its
        projected closing stock becomes the opening stock of the next period (Sections~\ref{sec:work-orders}
        to~\ref{sec:close}).
\end{enumerate}
